\makeabstract
    {
    Hypermethylation signature is associated with pretreatment resistance to Azacitidine-Venetoclax in acute myeloid leukemia
    }
    {
    Sofia I. Bianconi\textsuperscript{1}, Tamrin Chowdhury\textsuperscript{1}, Colton T. Hemphill\textsuperscript{1}, Rutu S. Vyas\textsuperscript{1}, Donald C. Green\textsuperscript{1}, Jack A. Gazil\textsuperscript{1}, Patrick W. Mellors\textsuperscript{2}, Parth S. Shah\textsuperscript{1,2}
    }
    {
    \textsuperscript{1} Department of Pathology and Laboratory Medicine, Dartmouth-Hitchcock Medical Center, Lebanon, NH\\
\textsuperscript{2} Department of Hematology / Oncology, Dartmouth-Hitchcock Medical Center, Lebanon, NH
    }
    {
    Acute myeloid leukemia (AML) is the most common type of leukemia affecting adult patients. In the last decade, Azacitadine and Venetoclax combination (Aza-Ven) became the standard of care treatment for elder AML patients unfit for more intensive chemotherapies. Despite its clinical efficacy, a subset of patients exhibits primary refractory disease, and the biological mechanisms underlying Aza-Ven resistance remain incompletely understood. In this study we investigated differential methylation between Aza-Ven sensitive and refractory bone marrow samples to identify a candidate epigenetic signature associated with refractory disease and explore its potential utility for pretreatment response stratification.
    }
    {
    Pretreatment bone marrow samples from 23 AML patients treated with Aza-Ven were collected and profiled using Oxford Nanopore long-read sequencing. Methylation pileup files were generated using modkit and differential methylated regions were identified using both 1kb genomic tiles and CpG island-based analyses. Comparisons were performed in the full Aza-Ven cohort, comprising 14 sensitive and 9 refractory patients, and in a blast-enriched subgroup restricted to samples with \ensuremath{\geq}40\% marrow blasts, comprising 12 sensitive and 7 refractory patients, to reduce confounding from non-leukemic cellular composition. Candidate regions were annotated to promoters, gene bodies, and nearby genes. Regions showing at least a 0.15 absolute difference in mean methylation between response groups, corresponding to 15 percentage points, and an adjusted P value \ensuremath{<}0.05 were prioritized. A candidate refractory-associated hypermethylation signature was defined from loci that met either one or both of these criteria and showed additional support through persistence in the blast-enriched analysis, concordance across tile and CpG-island resolutions. The potential predictive relevance of the signature was evaluated based on its presence in pretreatment samples and the consistency of higher methylation among refractory patients.
    }
    {
    A candidate hypermethylation signature associated with Aza-Ven refractory AML was identified and included loci annotated to NRGN, ARHGAP35, TNRC6C, HOXB4, HOXA6, HOXB3, SLC2A14, and LINC01347. The mean methylation difference between response groups of these loci ranged from 0.15 to 0.26 (15-26\% increase) with an adjusted $P$ value of $0.05$. NRGN provided the strongest cross-resolution evidence, demonstrating gene-body hypermethylation in refractory samples in both 1kb tile (21\% increase, P \ensuremath{<}0.007) and CpG island-based (20\% increase, P \ensuremath{<}0.02) analyses. Additional components of the signature were identified primarily in the blast-enriched analysis, suggesting that restriction to samples with greater leukemic-cell content improved detection of leukemia-associated methylation differences. Several signature genes, including NRGN, ARHGAP35, HOXA6, HOXB3, and HOXB4, are associated with hematopoietic developmental or stem/progenitor-related biological programs. Because these methylation differences were detected before treatment and were consistently enriched in refractory samples, the combined signature may have value for identifying patients with an epigenetically distinct leukemia state associated with reduced Aza-Ven responsiveness.
    }
    {
    Pretreatment Aza-Ven refractory AML was associated with a candidate hypermethylation signature involving genes related to hematopoietic developmental regulation, stem/progenitor-associated biology, cellular signaling, and transcriptional control. Cross-resolution and blast-enriched analyses provided complementary support for this signature. These findings suggest that pretreatment methylation profiling may help distinguish patients at increased risk of refractory disease. Independent validation and development of a sample level methylation score with formal classification and cross-validation will be required to determine the predictive accuracy and clinical utility of this signature.
    }

    \newpage
    